\documentclass[../../../main.tex]{subfiles}
\begin{document}

\section{Introduction: How to Read This Book}

We wrote this book with the lofty goal to aggregate everything we know in mathematics and the adjacent fields like computer science and statistics. Our own backgrounds are in mathematics, philosophy, statistics, and computer science, and the book will undoubtedly reflect those backgrounds. While we include some topics --- such as differential equations --- because we believe they are important to a complete picture of mathematics despite intially knowing little about them, for the most part the inclusion of topics generally reflects our own particular educational pathways. We make no claims about what topics are most ``natural" or most important. 

Therefore, we expect nobody to read this book from back to front. Instead, each chapter can generally be treated as standalone; if you want to learn about group theory, you can just start in the group theory section. Why, then, organize such a massive book, rather than a large number of smaller books? One benefit is that we try to enforce some common notation and ideas throughout the book, giving readers the opportunity to find any piece of prerequisite information they need without having to sift through a pile of resources that all use different notations, assume different prerequisites, and generally have different audiences in mind. (This is the bane of our existence.) Indeed, we hope that a reader could \textit{in principle} learn about any subject in this book without needing another resource, albeit with a lot of time and patience. 

Unfortunately, we do have to have some minimal prerequisites, as any book must. We hope that the prerequisites here are minimal. We start with  chapters on sets and logic, linear algebra, and calculus, which establish the foundations for the rest of the book. Of these, the calculus chapter moves somewhat quickly, assuming some prior exposure; otherwise, only basic mathematical skills (manipulating algebraic equations, connecting equations to graphical representations, careful and logical reasoning) are necessary. We hope that any high school student who covered up to a calculus class equivalent to AP Calculus BC will find the book approachable. 

How should you read this book? Start with a plan. What subject do you want to build up to? In the relevant chapter, we list the prerequisites in the preamble. If you feel comfortable with those prerequisites, you can start! If not, you may want to read the corresponding prerequisite chapters first, though if you're in a hurry you're welcome to try anyway and go back whenever you get stuck. As with any math textbook, you'll want to work slowly with pen and paper in hand, being sure to understand the definitions and key results thoroughly before progressing. 

What does this book cover?
Each ``Chapter" of the book roughly corresponds to one university course's worth of material, although this will of course vary. The Chapters are organized into ``Parts".
The first part --- ``Foundations" --- covers three subjects that are prerequisites for many mathematical topics: calculus, linear algebra, and mathematical foundations (logic, sets, and proofs). If you don't feel familiar with these subjects, you should start here; this is material that almost every math students has taken by their first year of college. The next two parts --- ``Introduction to Pure Math", and ``Introduction to Applied Math" --- cover material that pure and applied math students (respectively) would likely cover by years 2--3 of an undergraduate degree. These parts cover foundational subjects that are largely separate from each other, so any of these chapters can be read on their own. 

The following parts are oriented around subjects --- analysis, algebra, probability and statistics, etc. --- grouped somewhat subjectively based on topics we think belong together. (Does machine learning belong to statistics or computer science? Yes.) Some of these parts are written to be more linear, with chapters building upon one another, while others are better thought of as collections of standalone chapters. Each chapter is organized into sections, which roughly correspond to the length of one to a few lectures.

\end{document}
